% The Glashow resonance: an electron antineutrino hits an electron in the ice and makes a
% real W- (s-channel), at E = M_W^2 / 2 m_e ~ 6.3 PeV. The W then decays; drawn here to quarks
% (W- -> d ubar, ~2/3 of decays), as in IceCube's 2016 event (Nature 591, 220 (2021)).
% The other ~1/3: W- -> e- nubar_e, mu- nubar_mu (a track), tau- nubar_tau.
% Time runs upwards, as in nu-e-reactions.tex (same layout as its annihilation diagram, Z -> W).
% Uses feynmp via feynmp-auto.
% Build (in this directory):
%   pdflatex -shell-escape glashow.tex && pdflatex -shell-escape glashow.tex
%   pdftocairo -svg glashow.pdf ../glashow.svg
\documentclass[border=6pt]{standalone}
\usepackage{feynmp-auto}
\usepackage{xcolor}

\newif\ifdark \darktrue          % \darkfalse = black on white
\newcommand{\figfont}{\LARGE}    % font size for labels and captions

\ifdark \pagecolor{black} \fi
\unitlength=1mm

\begin{document}
\ifdark \color{white} \fi
\figfont
\begin{fmffile}{glashowW}
  \ifdark \fmfcmd{foreground := white; background := black;} \fi
  \fmfset{thin}{0.8pt}
  \fmfset{arrow_len}{3mm}
  \fmfset{wiggly_len}{3mm}
  \begin{tabular}{c}
    $\vcenter{\hbox{\fmfframe(4,8)(4,8){% margins (left,top)(right,bottom) leave room for labels
    \begin{fmfgraph*}(28,40)
      \fmfbottom{nub,em}
      \fmftop{d,ub}
      \fmf{fermion}{em,v1,nub}
      \fmf{boson,label=$W^-$,label.side=right}{v1,v2}
      \fmf{fermion}{ub,v2,d}
      \fmflabel{$\bar\nu_e$}{nub}
      \fmflabel{$e^-$}{em}
      \fmflabel{$d$}{d}
      \fmflabel{$\bar u$}{ub}
    \end{fmfgraph*}}}}$
    \\[2mm]
    $\bar\nu_e\, e^- \to W^- \to d\,\bar u$
  \end{tabular}
\end{fmffile}
\end{document}
